\section{Introduction}\label{sec:s_1}

\subsection{Motivations}\label{subsec:ss1_1}


\subsection{Etat de l'art}\label{subsec:ss1_2}


\section{Quelques preuves}\label{sec:s_2}

\subsection{Le cas des graphes non connexes}\label{subsec:ss2_1}

On montre dans cette partie que tout graphe non connexe admet une partition satisfaisante.

Soit \(G = (V, E)\) un graphe non connexe.
Soit \(A\) une composante connexe de \(G\) et \(B = V \setminus A\) le reste des sommets.

On considère la partition \((A, B)\) de \(V\).
Soit \(x\) un sommet de \(A\). Le degré interne de \(x\) dans \(A\) est au moins 1 car \(A\) est une composante connexe.
Son degré externe (donc dans \(B\)) est de 0 car il n'y a pas d'arêtes entre \(A\) et \(B\).
Ainsi, pour tout sommet \(x\) de \(A\), on a \(d_A(x) \geq d_B(x)\).

Soit \(y\) un sommet de \(B\). Le degré externe de \(y\) dans \(A\) est de 0 car il n'y a pas d'arêtes entre \(A\) et \(B\).
Comme le degré ne peut pas être négatif, on a \(d_B(y) \geq 0\).
Ainsi, pour tout sommet \(y\) de \(B\), on a \(d_B(y) \geq d_A(y)\).

On a donc montré que pour tout sommet \(x\) de \(V\), le degré interne de \(x\) est supérieur ou égal à son degré externe pour la partition \((A, B)\).
Ainsi, \((A, B)\) est une partition satisfaisante de \(G\).

Dans la suite on ne considérera donc que des graphes connexes, car les graphes non connexes admettent toujours une partition satisfaisante.

\subsection{Les cas \(k \leq 4\)}\label{subsec:ss2_2}

\subsubsection{Le cas \(k = 1\)}\label{subsubsec:ss2_2_1}

Soit \(G\) un graphe \(1\)-régulier connexe.
\(G\) est forcément le graphe à deux sommets reliés par une arête.
Il n'admet pas de partition satisfaisante, car les deux sommets ont un degré interne de 0 et un degré externe de 1.

\subsubsection{Le cas \(k = 2\)}\label{subsubsec:ss2_2_2}

Soit \(G\) un graphe \(2\)-régulier connexe.
\(G\) est un cycle de longueur \(n \geq 3\).
Soit \(x\) un sommet de \(G\) et \(y\) un sommet adjacent à \(x\).

Si \(n = 3\), alors \(G\) est un triangle et n'admet pas de partition satisfaisante car le sommet dans la partie de taille 1 a un degré interne de 0 et un degré externe de 2.

Si \(n \geq 4\), alors on peut partitionner les sommets de \(G\) en deux parties: \(({x, y}, V \setminus \{x, y\})\).

Soit \(z\) un sommet de \(G\).
Si \(z\) est dans la partie \(\{x, y\}\), alors son degré interne est de 1 (car il est adjacent à l'autre sommet de la partie) et son degré externe est de 1 (car il est adjacent à un sommet de l'autre partie).
Si \(z\) est dans la partie \(V \setminus \{x, y\}\), alors son degré interne est de au moins 1 (il ne peut pas être lié à \(x\) et \(y\) car \(n \geq 4\)) et son degré externe est de au plus 1.

Ainsi, pour tout sommet \(z\) de \(G\), on a \(d_{\mathrm{int}}(z) \geq d_{\mathrm{ext}}(z)\) pour la partition \(({x, y}, V \setminus \{x, y\})\).
Donc \(G\) admet une partition satisfaisante.

\subsubsection{Le cas \(k = 3\)}

Soit \(G = (V,E)\) un graphe \(3\)-régulier connexe.
On va montrer que \(G\) admet une partition satisfaisante, sauf si \(G\) est isomorphe à \(K_4\) ou \(K_{3,3}\).

Soit \((A,B)\) une partition de \(V\) telle que \(A \neq \emptyset\) et \(B \neq \emptyset\).

On dit qu’un sommet \(v \in V\) est \emph{mauvais} si \(d_{\mathrm{ext}}(v) > d_{\mathrm{int}}(v)\).
Comme \(G\) est \(3\)-régulier, cela équivaut à \(d_{\mathrm{ext}}(v) \geq 2\).

On considère le processus suivant: tant qu’il existe un sommet mauvais \(v\), on le déplace dans l’autre partie, à condition que sa partie contienne au moins deux sommets.

On définit la fonction \(f\) comme suit:
\begin{equation*}
    f(A,B) = |E(A)| + |E(B)|
\end{equation*}
où \(E(A)\) (resp. \(E(B)\)) désigne l’ensemble des arêtes internes à \(A\) (resp. \(B\)).

On considère alors les propositions suivantes:

\begin{proposition}
    A chaque déplacement autorisé, la fonction \(f\) augmente strictement.
\end{proposition}

\begin{proof}
    Soit \(v\) un sommet mauvais. Alors \(d_{\mathrm{ext}}(v) \geq 2\) et donc \(d_{\mathrm{int}}(v) \leq 1\).
    En déplaçant \(v\), on perd \(d_{\mathrm{int}}(v)\) arêtes internes et on gagne \(d_{\mathrm{ext}}(v)\) arêtes internes.
    Ainsi,
    \[
    \Delta f \geq 2 - 1 = 1.
    \]
    Donc \(f\) augmente strictement.
\end{proof}

Comme \(f\) est bornée par \(|E|\), le processus termine.

\begin{proposition}\label{prop:2}
    Si le processus s’arrête et que la partition n’est pas satisfaisante, alors l’une des deux parties est un singleton.
\end{proposition}

\begin{proof}
    Si un sommet mauvais appartient à une partie de taille au moins 2, alors il pourrait être déplacé, ce qui contredit l’arrêt du processus.
    Donc tout sommet mauvais appartient à une partie de taille 1.
\end{proof}

Supposons maintenant que le processus s’arrête sur une partition non satisfaisante. Par la proposition~\ref{prop:2}, on peut supposer sans perte de généralité que \(A = \{v\}\). % chktex 44
Alors \(v\) a ses trois voisins dans \(B\), qu'on note \(x\), \(y\) et \(z\).
On analyse alors la structure de \(B\). 

Si \(B\) contient exactement \(3\) sommets alors \(G\) est isomorphe à \(K_4\).

\(B\) ne peut pas contenir exactement \(4\) sommets car \(G\) est \(3\)-régulier.

Si \(B\) contient exactement \(5\) sommets on a 2 cas possibles:
\begin{itemize}
    \item soit \(x\), \(y\) et \(z\) sont tous les trois liés à un même sommet de \(B\), alors \(G\) est isomorphe à \(K_{3,3}\),
    \item soit \(x\), \(y\) et \(z\) ne sont pas tous les trois liés à un même sommet de \(B\), alors \(G\) contient deux cycles disjoints de longueur 3 et on a trouvé une autre partition satisfaisante.
\end{itemize}

% to fill
\medskip

Ainsi, si \(G\) n’est isomorphe ni à \(K_4\) ni à \(K_{3,3}\), le processus ne peut pas s’arrêter sur une partition non satisfaisante.

On en déduit que \(G\) admet une partition satisfaisante.